% Edit and compile this template online at https://letx.app
\documentclass[10pt,letterpaper]{article}

\usepackage[utf8]{inputenc}
\usepackage[T1]{fontenc}

% Use Lato font as specified in AGY_BRIEF.md
\usepackage{lato}
\renewcommand{\familydefault}{\sfdefault}

% Margins and header spacing
\usepackage[margin=0.75in, top=0.95in, bottom=0.95in, headheight=15pt]{geometry}

% Core packages
\usepackage{xcolor}
\usepackage{tabularx}
\usepackage{booktabs}
\usepackage{enumitem}
\usepackage{titlesec}
\usepackage[most]{tcolorbox}
\usepackage{amsmath}
\usepackage{siunitx}
\usepackage{microtype}
\usepackage{fancyhdr}
\usepackage{hyperref}

% Define tasteful scientific palette (deep teal/slate blue)
\definecolor{accent}{HTML}{0D6E6E}       % Deep Teal
\definecolor{accentlight}{HTML}{F2F7F7}  % Light Teal Tint
\definecolor{accentdark}{HTML}{073B3B}   % Dark Teal
\definecolor{neutraldark}{HTML}{1E293B}  % Slate Dark (text color)
\definecolor{neutralgray}{HTML}{64748B}  % Slate Gray
\definecolor{bordercolor}{HTML}{CBD5E1}  % Slate Light Border
\definecolor{bodyback}{HTML}{F8FAFC}     % Slate Page Tint

% Set text color globally
\color{neutraldark}

% Disable indent and parskip inside standard formatting to align tables and prevent overflow
\setlength{\parindent}{0pt}
\setlength{\parskip}{0pt}

% Configure lists with enumitem
\setlist[enumerate]{leftmargin=18pt, labelsep=6pt, itemsep=4pt, parsep=0pt, topsep=4pt}
\setlist[itemize]{leftmargin=18pt, labelsep=6pt, itemsep=4pt, parsep=0pt, topsep=4pt}

% Hyperref setup
\hypersetup{
  colorlinks=true,
  linkcolor=accent,
  urlcolor=accent,
  pdfauthor={LetX},
  pdftitle={Elegant Lab Notebook Template}
}

% Page styling with fancyhdr
\pagestyle{fancy}
\fancyhf{}
\renewcommand{\headrulewidth}{0pt}
\renewcommand{\footrulewidth}{0.4pt}
\fancyfoot[L]{\color{neutralgray}\footnotesize \textbf{CONFIDENTIAL} -- Internal Research Log}
\fancyfoot[R]{\color{neutralgray}\footnotesize Page \thepage}

% Custom section command for structured entry fields
\newcommand{\labsection}[1]{%
  \vspace{10pt}
  {\color{accent}\noindent\large\bfseries\sffamily #1}
  \vspace{3pt}
  {\color{bordercolor}\hrule height 0.8pt}
  \vspace{5pt}
}

% Custom header block command
\newcommand{\labheader}[4]{%
  \begin{tcolorbox}[
    enhanced,
    colback=accentlight,
    colframe=accent,
    arc=4pt,
    outer arc=4pt,
    boxrule=1.5pt,
    left=12pt, right=12pt, top=10pt, bottom=10pt,
    boxsep=0pt,
    shadow={2pt}{-2pt}{0pt}{accent!10!white}
  ]
    \begin{tabularx}{\linewidth}{@{}>{\color{accent}\bfseries}l@{\hspace{8mm}}>{\raggedright\arraybackslash}X@{\hspace{12mm}}>{\color{accent}\bfseries}l@{\hspace{8mm}}>{\raggedright\arraybackslash}X@{}}
      Project: & #1 & Date: & #3 \\
      Researcher: & #2 & Entry No: & #4 \\
    \end{tabularx}
  \end{tcolorbox}
  \vspace{4pt}
}

% Ruled/lightly-tinted entry area environment
\newtcolorbox{entryarea}{
  enhanced,
  boxrule=0.5pt,
  colframe=bordercolor,
  colback=bodyback,
  arc=4pt,
  left=15pt, right=15pt, top=10pt, bottom=15pt,
  grow to left by=2pt, grow to right by=2pt,
  boxsep=0pt,
  before skip=5pt, after skip=5pt
}

\begin{document}

% ==================== PAGE 1 (Day 1 Entry) ====================
\labheader{Nanomedicine Synthesis}{Dr. Evelyn Vance}{July 5, 2026}{104-A}

\begin{entryarea}
  \labsection{Objective}
  To synthesize spherical silver nanoparticles (AgNPs) via chemical reduction of silver nitrate using sodium citrate as both the reducing agent and stabilizing capping agent. Size target is $20\,\text{nm}$.

  \labsection{Materials \& Reagents}
  \begin{tabularx}{\linewidth}{@{}>{\color{accent}\bfseries}l@{\hspace{8mm}}>{\raggedright\arraybackslash}X@{}}
    Silver Nitrate ($\text{AgNO}_3$): & $1.0\,\text{mM}$ aqueous stock solution, 99.9\% purity (Sigma-Aldrich). \\
    Sodium Citrate: & $1\%$ (w/v) tribasic dihydrate aqueous solution (Analytical Grade). \\
    Solvent: & Ultra-pure deionized water ($18.2\,\text{M}\Omega\cdot\text{cm}$ resistivity). \\
    Apparatus: & Cleaned 250 mL Pyrex Erlenmeyer flask, magnetic stirrer, thermometer. \\
  \end{tabularx}

  \labsection{Procedure}
  \begin{enumerate}
    \item Clean all glassware with aqua regia followed by thorough rinsing with deionized water to prevent premature seeding and contamination.
    \item Prepare $100\,\text{mL}$ of $1.0\,\text{mM}$ aqueous solution of silver nitrate ($\text{AgNO}_3$) in the 250 mL Erlenmeyer flask.
    \item Heat the solution on a magnetic stirrer hotplate to a vigorous boil under constant stirring ($600\,\text{rpm}$).
    \item Quickly add $10\,\text{mL}$ of $1\%$ (w/v) sodium citrate solution to the boiling mixture.
    \item Maintain boiling temperature and constant stirring for 25 minutes. Note the color changes at 5-minute intervals.
    \item Remove from heat and allow the colloidal suspension to cool to room temperature. Store in a dark amber vial.
  \end{enumerate}

  \labsection{Observations}
  \begin{itemize}
    \item \textbf{0 min:} Initial solution is completely colorless and clear.
    \item \textbf{5 min:} A faint yellow hue begins to develop, indicating nucleation.
    \item \textbf{12 min:} Color intensifies to a bright, golden yellow, indicating excitation of the surface plasmon resonance (SPR).
    \item \textbf{25 min:} Stable deep yellow/amber colloidal suspension. No precipitate or silver mirroring on the flask walls.
  \end{itemize}
\end{entryarea}

\newpage

% ==================== PAGE 2 (Day 2 Entry) ====================
\labheader{Nanomedicine Characterization}{Dr. Evelyn Vance}{July 6, 2026}{104-B}

\begin{entryarea}
  \labsection{Objective}
  To characterize the synthesized silver nanoparticles from Entry 104-A using UV-Vis spectrophotometry to determine peak absorption ($\lambda_{\max}$) and evaluate size distribution and monodispersity.

  \labsection{Materials \& Reagents}
  \begin{tabularx}{\linewidth}{@{}>{\color{accent}\bfseries}l@{\hspace{8mm}}>{\raggedright\arraybackslash}X@{}}
    Colloidal AgNPs: & Synthesized solution from Entry 104-A (stored at $4\,^{\circ}\text{C}$). \\
    Blank Solvent: & Deionized water ($18.2\,\text{M}\Omega\cdot\text{cm}$ resistivity). \\
    Apparatus: & UV-Vis Spectrophotometer (Shimadzu UV-1900i), 10 mm quartz cuvettes. \\
  \end{tabularx}

  \labsection{Procedure}
  \begin{enumerate}
    \item Turn on the spectrophotometer and warm up the light source for 15 minutes.
    \item Perform baseline calibration using deionized water in both reference and sample cuvettes from $350\,\text{nm}$ to $700\,\text{nm}$.
    \item Dilute the synthesized AgNP colloidal suspension 1:5 with deionized water to prevent detector saturation.
    \item Rinse the sample cuvette with a small aliquot of the diluted suspension, then fill.
    \item Run the wavelength scan and record the absorbance spectra.
    \item Identify the wavelength corresponding to peak absorbance ($\lambda_{\max}$) and calculate the full width at half maximum (FWHM).
  \end{enumerate}

  \labsection{Results \& Data Analysis}
  A summary of the spectrophotometric analysis for three consecutive dilution runs is tabulated below:
  
  \vspace{6pt}
  \begin{center}
    \small
    \begin{tabularx}{\linewidth}{@{}>{\raggedright\arraybackslash}X c c c c@{}}
      \toprule
      \textbf{Sample Run} & \textbf{Peak $\lambda_{\max}$ (nm)} & \textbf{Absorbance (a.u.)} & \textbf{FWHM (nm)} & \textbf{Est. Size (nm)} \\
      \midrule
      AgNP-104A-Run1 & 418.5 & 0.842 & 54.2 & 18.5 \\
      AgNP-104A-Run2 & 420.1 & 0.865 & 55.0 & 20.2 \\
      AgNP-104A-Run3 & 419.8 & 0.859 & 54.8 & 19.8 \\
      \midrule
      \textbf{Mean ($\pm$ SD)} & \textbf{419.5 $\pm$ 0.8} & \textbf{0.855 $\pm$ 0.012} & \textbf{54.7 $\pm$ 0.4} & \textbf{19.5 $\pm$ 0.9} \\
      \bottomrule
    \end{tabularx}
  \end{center}
  \vspace{2pt}

  \labsection{Conclusion \& Next Steps}
  The UV-Vis spectrum shows a sharp, symmetrical SPR peak at $419.5\,\text{nm}$, indicating the successful synthesis of spherical silver nanoparticles with an estimated diameter of $\sim 20\,\text{nm}$. The low FWHM ($54.7\,\text{nm}$) confirms a narrow size distribution and good monodispersity.
  
  \textbf{Next steps:} Perform Transmission Electron Microscopy (TEM) and Dynamic Light Scattering (DLS) to confirm nanoparticle morphology and hydrodynamic size.
\end{entryarea}

\end{document}
